\section{Introduction}
\IEEEPARstart{S}{upervisory} recipe adjustment links a production objective to
changes on the line. Film thickness, for example, depends on screw speed,
line speed, and temperatures; biaxial stretching couples adjustments across
several sections. Choosing a setpoint is only one part of this task. Each
intervention also needs execution authority, applicable operating limits, and
observations that distinguish a completed command from an improved process.

Large language model (LLM) agents are being explored for industrial diagnosis, engineering
assistance, and programmable logic controller (PLC) program generation
\cite{raza2025industrial,sabetta2025assistant,liu2026agents4plc,fakih2024llm4plc},
including interaction with simulated controllers and machine demonstrators
\cite{vyas2026ftc,hofmann2025nlc}. As agents coordinate interventions across
equipment, protocol connectivity must be coupled to supervisory context:
which deployed member may change which parameter, under which recipe, and
with whose approval? Task orchestration and device access alone do not
answer these questions or connect a write to its subsequent judgment.

\system{} makes these relationships explicit through member--node bindings
and shared context for lines, products, recipes, and runs. Nodes expose engineering
quantities rather than register addresses; bindings grant scoped tool
authority; and intervention records associate changes with observations and
submitted judgments. Fig.~\ref{fig:intro} contrasts this organization with
application-specific integration. The aim is to make supervisory search
inspectable across heterogeneous protocols while keeping fast regulation and
plant protection outside the agent loop.

\begin{figure*}[!tb]
\centering
\includegraphics[width=\textwidth]{figures/publication/fig1-supervisory-binding.pdf}
\caption{Conceptual supervisory integration: (a)~conventional
application-specific connections; (b)~explicit member--node capabilities and
shared production context in the proposed organization. PLC regulation remains
outside the agent loop. The panels illustrate design patterns, not an
empirical comparison of competing systems.}
\label{fig:intro}
\end{figure*}

At seconds-to-minutes timescales, manual bindings require approval and
automatic bindings retain execution checks without that wait. Separate
judgment and compensation operations distinguish rejection, driver acceptance,
readback agreement where supported, process response, and attempted
restoration. These distinctions define three contributions:
\begin{enumerate}
\item A \textbf{node-native supervisory model} connecting heterogeneous
acquisition and control nodes, deployed members, and production context through
explicit capabilities and shared identifiers
(Secs.~\ref{sec:arch} and \ref{sec:objective}).
\item An \textbf{implemented governed loop} linking source-dependent admission,
intervention records with optional hypotheses, submitted judgments, separate
compensation, and human approval (Sec.~\ref{sec:mech}).
\item A \textbf{layered, traceable evaluation} using protocol checks, an
interlock ablation, fixed-controller trials, and scripted team missions to
characterize constraint enforcement and multi-node loop closure
(Sec.~\ref{sec:eval}).
\end{enumerate}
The evaluation uses real protocol transports with simulated plants and
deterministic policies, isolating supervisory execution from LLM reasoning.
The contribution is an execution and evidence framework, not a new control
law or a claim of field superiority. Configured recipe windows constrain
proposals; neither they nor shared context provide atomic execution or a
physical safety guarantee.
